\documentclass[a4paper,11pt]{article}
\usepackage[a4paper,margin=2cm]{geometry}
\usepackage[T1]{fontenc}
\usepackage[utf8]{inputenc}
\usepackage[ngerman,english]{babel}
\usepackage{lmodern}
\usepackage{microtype}
\usepackage{xcolor}
\usepackage{parskip}
\usepackage{enumitem}
\usepackage[most]{tcolorbox}
\usepackage{titlesec}
\usepackage[hidelinks]{hyperref}
\usepackage{array}
\usepackage{longtable}
\usepackage[table]{xcolor}
\definecolor{HeaderBlueDark}{RGB}{70,110,170}
\definecolor{RowBlueLight}{RGB}{235,243,255}
\definecolor{RowBlueDark}{RGB}{220,233,250}
\definecolor{SoftGray}{RGB}{90,90,90}
\setlength{\parindent}{0pt}
\setlist[itemize]{leftmargin=1.4em,itemsep=0.35em,topsep=0.35em}
\linespread{1.06}
\renewcommand{\arraystretch}{1.35}
\setlength{\tabcolsep}{5pt}
\titleformat{\section}[block]{\Large\bfseries\color{HeaderBlueDark}}{}{0pt}{}
\titlespacing{\section}{0pt}{1.3em}{0.8em}
\titleformat{\subsection}[block]{\large\bfseries\color{HeaderBlueDark}}{}{0pt}{}
\titlespacing{\subsection}{0pt}{1.0em}{0.6em}
\newtcolorbox{exbox}{enhanced,breakable,colback=RowBlueLight,colframe=RowBlueDark,boxrule=0.4pt,arc=1.5mm,left=1.2mm,right=1.2mm,top=1mm,bottom=1mm,title={Beispiele \textnormal{(Examples)}},fonttitle=\bfseries,colbacktitle=RowBlueDark,coltitle=black}
\NewDocumentEnvironment{grammarentry}{mm+b}{%
\begin{tcolorbox}[enhanced,breakable,colback=white,colframe=HeaderBlueDark,boxrule=0.6pt,arc=1.5mm,left=1.5mm,right=1.5mm,top=1mm,bottom=1mm,colbacktitle=HeaderBlueDark,coltitle=white,fonttitle=\bfseries,title={#1\hfill\normalfont\footnotesize #2}]
#3
\end{tcolorbox}}{}
\newcommand{\GermanText}[1]{\textbf{Deutsch: }#1\par}
\newcommand{\EnglishText}[1]{{\small\color{SoftGray}\textbf{English: }#1\par}}
\newcommand{\ExampleItem}[2]{\item \textbf{#1}\\{\small\color{SoftGray}#2}}
\newcommand{\DEEN}[2]{\textbf{#1}\par{\footnotesize\color{SoftGray}#2}}
\title{Passiv im Konjunktiv\\[0.3em]\large\textnormal{Passive in the subjunctive}}
\date{}
\begin{document}
\maketitle
\tableofcontents
\newpage

\section{Konjunktiv I -- Subjunctive I}
\begin{grammarentry}{Indirekte Rede}{Reported speech}
\GermanText{Der Konjunktiv I wird besonders in der indirekten Rede benutzt. Im Vorgangspassiv steht häufig \textbf{werde + Partizip II}; im Zustandspassiv \textbf{sei + Partizip II}.}
\EnglishText{Subjunctive I is used especially in reported speech. In the event passive, \textbf{werde + past participle} is common; in the state passive, \textbf{sei + past participle}.}
\end{grammarentry}
\begin{exbox}\begin{itemize}
\ExampleItem{Er sagt, die Arbeit werde morgen erledigt.}{He says that the work will/is said to be done tomorrow.}
\ExampleItem{Er sagt, die Arbeit sei schon erledigt.}{He says that the work is already completed.}
\end{itemize}\end{exbox}

\section{Konjunktiv II -- Subjunctive II}
\begin{grammarentry}{Irreales und Hypothetisches}{Unreal and hypothetical meaning}
\GermanText{Im Vorgangspassiv wird häufig \textbf{würde + Partizip II + werden} verwendet. Im Zustandspassiv steht häufig \textbf{wäre + Partizip II}.}
\EnglishText{In the event passive, \textbf{würde + past participle + werden} is commonly used. In the state passive, \textbf{wäre + past participle} is common.}
\end{grammarentry}
\begin{exbox}\begin{itemize}
\ExampleItem{Das Problem würde gelöst werden.}{The problem would be solved.}
\ExampleItem{Das Problem wäre gelöst.}{The problem would be solved / would be in a solved state.}
\end{itemize}\end{exbox}

\section{Mit Modalverben -- With modal verbs}
\begin{grammarentry}{Konjunktiv II}{Subjunctive II}
\GermanText{Auch Modalverben können im Konjunktiv II mit dem Passiv stehen: \textbf{Die Aufgabe müsste heute erledigt werden.}}
\EnglishText{Modal verbs can also be used in Subjunctive II with the passive: \textbf{The task would have to be completed today.}}
\end{grammarentry}

\section{Merksatz -- Key rule}
\begin{grammarentry}{Kurz}{In short}
\GermanText{Konjunktiv verändert die Modalität, während die Passivstruktur erhalten bleibt.}
\EnglishText{The subjunctive changes the modality while the passive structure remains intact.}
\end{grammarentry}
\end{document}
